\section{A Compact Mean-Width Specialization}
\label{sec:structured}\label{sec:randomized-recovery}
The general hardest-record LP can enumerate exponentially many batches.
A supporting specialization avoids this enumeration for disjoint exact
totals and coordinate queries. It has a different measurement interface
from cumulative-register retrieval: each query returns one missing increment.
Write $\F=\prod_g\{y_g\geq0:\sum_i y_{gi}=E_g\}$ and let $d$ include
all known contributions to the affine contrast. Its sharp endpoints are
\begin{equation}
 D^-=d+\sum_gE_g\min_i c_{gi},\qquad
 D^+=d+\sum_gE_g\max_i c_{gi}.\label{eq:sharp}
\end{equation}
For retained coordinates $S_g(Q)$, with singleton/empty ranges defined as zero,
\begin{equation}
 W(Q,y)=\sum_g\left(\sum_{i\in S_g(Q)}y_{gi}\right)
             \operatorname{range}(c_{g,S_g(Q)}).
 \label{eq:querywidth}
\end{equation}
Define $v_{gi}(Q)=\mathbf1\{i\in S_g(Q)\}\operatorname{range}(c_{g,S_g(Q)})$.
Omit zero-width blocks and redundant queries of a last determined coordinate.
For admissible batches $\mathcal Q_B$ with at most $B$ unit-cost queries,
\begin{equation}
 W(Q,y)=\sum_{g,i}y_{gi}v_{gi}(Q),\qquad
 R_B=\min_{\pi\in\Delta(\mathcal Q_B)}\sum_gE_g\max_i\mathbb E_\pi v_{gi}(Q).
 \label{eq:random-objective}
\end{equation}
Maximizing a linear function on each mass simplex proves this expression.
The resulting finite LP has primal constraints
$\mathbb E_\pi v_{gi}(Q)\leq t_g$ and objective $\min\sum_gE_gt_g$;
its dual is
\begin{equation}
 \max_{w,\eta}\eta\quad\text{subject to}\quad
 w_{gi}\geq0,\quad\sum_iw_{gi}=E_g,\quad
 \eta\leq\sum_{g,i}w_{gi}v_{gi}(Q)\quad(Q\in\mathcal Q_B).
 \label{eq:random-dual}
\end{equation}

The separation oracle sorts coefficients, fixes the two retained endpoints
and selects the smallest interior weights until the retained cardinality is
met. A budget-allocation dynamic program then joins the blocks. Weighted
interiors need not form consecutive coefficient windows. The equivalent
compact LP represents this dynamic program by a unit flow through block/budget
states. Each edge fixes retained endpoints; conditional interior-inclusion
variables satisfy $0\leq z_{ei}\leq p_e$ and
$\sum_i z_{ei}=(r_e-2)p_e$, where $p_e$ is edge flow and $r_e$ is retained
cardinality. These are scaled fixed-cardinality subset polytopes. Flow-path
decomposition and their integral vertices produce legal batch distributions;
conversely every legal batch induces such flows and marginals. Thus the
representation has no relaxation gap. For $B\geq1$, a safe unpruned size
bound is $O(B\sum_g\min(B,m_g)m_g^3)$, and $O(m_g^3)$ for one block with
fixed budget. At $B=0$, the sharp interval is evaluated directly from
\eqref{eq:sharp}, without constructing a recovery LP.
Online Resource~1, S5, provides the complete formulation, decomposition
proof, deterministic reference and weighted oracle. General weighted-oracle
robust optimization and dynamic-program representations are established
\citep{kawase2024,martin1990}; the retained-endpoint construction is the
specific implementation used here.

\begin{proposition}[A computable hardest completion]
\label{prop:adaptive-dual-completion}\label{sec:adaptive-dual-completion}
An optimal dual vector $w^*$ in \eqref{eq:random-dual}, with determined
and omitted independent blocks restored, is a feasible completion $y^*$.
Every admissible batch satisfies $W(Q,y^*)\geq R_B$. The compact batch
policy is therefore mean-width optimal even among randomized adaptive
coordinate-query policies with the same hard budget.
\end{proposition}
\begin{proof}
The nonnegative dual block masses are exactly the residual target
constraints. Strong LP duality gives $\eta^*=R_B$, so the dual batch
inequalities and \eqref{eq:random-objective} prove the stated bound.
For each fixed seed, the transcript-fiber argument of
Proposition~\ref{prop:risk_concavity_transcript} identifies the adaptive
terminal interval with that of its final queried set. At $y^*$ every
such execution has width at least $R_B$. The compact batch attains $R_B$
in worst expectation, giving the reverse inequality.
\end{proof}

The implementation's redundant nonnegative epigraph bounds can yield dual
block masses below $E_g$. Adding a deficit to a fixed coordinate cannot
decrease any payoff because $v_{gi}(Q)\geq0$; weak duality bounds the
completed payoff by the primal value. Numerical outputs retain raw masses,
tiny excess-mass corrections, and the completed record for pricing checks.
These are floating-point certificates with reported tolerances.

\paragraph{Choosing among equally robust policies.}\label{sec:tie-selection}
Following the optimal-face principle of \citet{iancu2014}, the secondary
objective minimizes $\sum_g(E_g/m_g)\sum_i\mathbb E_\pi v_{gi}(Q)$,
the width at the uniform-within-block completion, subject to attaining
$R_B$. It uses no hidden answers and removes pointwise-dominated expected
width functions. A separately checked refinement also requires never
worsening the archived policy. Neither refinement claims to optimize
precision or requests until a decision. Their adverse recorded outcomes
remain reported in the historical control study.

